\section{Overview}\label{4.1}

This chapter reports what was measured, in the order that best explains the
system's behaviour: how the benchmark is scored, what the raw model can do
unaided, where a pipeline genuinely has room to help, and how each added
component performs on its own terms.

\section{How Results Are Measured}\label{4.2}

Every answer in this project is scored deterministically against a
machine-checkable answer key, never by a second language model's
judgement (Section~\ref{2.5}). This section defines each metric exactly,
since the tables in the rest of the chapter are read against these
definitions.

\subsection*{Concept Recall}

Each benchmark question carries a set of required key phrases. A phrase
counts as present if it appears anywhere in the answer under a
normalisation that ignores case, markdown emphasis, British/American
spelling variants, and simple singular/plural agreement on the final word,
but performs no stemming and no synonym matching. For a question with
required phrases $H$, of which $h$ are matched:
\[
\text{concept recall} = \frac{|h|}{|H|}
\]

\subsection*{Avoided Forbidden}

A boolean flag: true if none of the question's forbidden phrases (common
misconceptions or incorrect claims) appear in the answer under the same
normalisation.

\subsection*{Numeric Correctness}

For questions with a declared numeric answer, the value stated in the
response is extracted and compared against the key within a stated
tolerance. A second, independent check re-derives the answer symbolically
from the question's own premises (Hardy--Weinberg, Mendelian ratios,
population growth, Michaelis--Menten, dilution, and carrier-probability
questions) and corroborates the extracted value rather than replacing it.

\subsection*{Order Correctness}

For questions describing a biological process, the answer's stated step
sequence is checked against a canonical ordering for that process.

\subsection*{Hallucination}

For the eight adversarial questions, which presuppose an entity that does
not exist, an answer is scored as hallucinating if it either states a
forbidden fact about the fictitious entity or fails to refuse the premise
at all.

\subsection*{Overall Score}

The four checks above are combined differently depending on the question
type, so that the score is always dominated by the check that matters most
for that kind of question:

\begin{description}
\item[Adversarial questions] score 1 if the answer correctly refuses the
false premise, 0 if it hallucinates.
\item[Numeric questions] score $0.8 \times \text{correct} + 0.2 \times
\text{concept recall}$.
\item[Process questions] score $0.5 \times \text{concept recall} + 0.4
\times \text{order correct} + 0.1 \times \text{avoided forbidden}$.
\item[All other questions] score $\max(0,\ \text{concept recall} - 0.5
\times \text{forbidden penalty})$, where the penalty applies only if a
forbidden phrase was stated.
\end{description}

\subsection*{Verifier Decision, Confidence, and Grounding}

Independently of benchmark scoring, every generated answer is passed
through the pipeline's own verification layer (Section~\ref{4.6}), which
reports three further quantities used throughout this chapter.

\begin{description}
\item[Decision] is \texttt{accept} unless a deterministic checker raises a
blocking error or a claim is materially contradicted by the evidence, in
which case it is \texttt{revise}, or unless essentially none of the
answer's claims tie to any retrieved passage, in which case it is
\texttt{insufficient\_evidence}.
\item[Grounded fraction] is the share of the answer's claims that overlap
sufficiently with at least one retrieved passage.
\item[Confidence] is the grounded fraction minus a penalty of $0.1$ per
blocking issue and $0.05$ per warning, clipped to $[0,1]$.
\end{description}

\subsection*{Provenance}

Three further numbers describe how traceable an answer is to its sources,
computed from the answer's inline citation markers and the passages
supplied to the model:

\begin{description}
\item[Claim citation rate] is the fraction of factual claims carrying at
least one inline citation.
\item[Citation validity] is the fraction of citations that point at a
passage plausibly containing the claim.
\item[Unsupported assertion rate] is $1$ minus the claim citation rate:
the share of claims with no citation at all.
\end{description}

\section{Baseline Capability and Available Headroom}\label{4.3}

Before any retrieval, verification, or critique is added, the raw local
model is measured on the full 441-question benchmark with no pipeline at
all: no retrieved evidence, no checking, nothing but the question and the
model.

\begin{table}[h]
\centering
\caption{Raw model performance, no pipeline, all 441 questions}
\label{tab:4.1}
\begin{tabular}{@{}lr@{}}
\toprule
\textbf{Metric} & \textbf{Score} \\
\midrule
Overall & 0.855 \\
Numeric accuracy & 0.800 \\
Process-order accuracy & 0.833 \\
\bottomrule
\end{tabular}
\end{table}

The headline result is that \textbf{55\% of the 441 questions are answered
perfectly by the unaided model}. This is the single fact that shapes most
of the decisions described in the rest of this chapter: a model that
already gets most of a benchmark right has very little room left for any
addition to demonstrably improve it.

To quantify that room directly, two independent measurements were made.
First, over the 50 questions used for the ablation study in
Section~\ref{4.11}, every required concept in the answer key was checked
against three things: whether it was present anywhere in the passages
retrieval returned, whether the unaided model named it, and their union.

\begin{table}[h]
\centering
\caption{Concept coverage ceiling, 50 questions, 188 required concepts}
\label{tab:4.2}
\begin{tabular}{@{}lr@{}}
\toprule
\textbf{Coverage of required concepts} & \textbf{Fraction} \\
\midrule
Present in retrieved evidence & 0.878 \\
Known by the unaided model & 0.878 \\
\textbf{Union of evidence and model memory} & \textbf{0.936} \\
\bottomrule
\end{tabular}
\end{table}

Retrieval and the model's own memory overlap on roughly 82\% of what the
benchmark asks for. Retrieval's unique contribution, the concepts the
model does not already know, is only the remaining 5.8\%. A perfect
pipeline could not beat the raw model here by more than about six points,
because the questions largely ask for exactly the kind of standard biology
knowledge a model this size has already absorbed during pretraining.

Second, a separate 38-question set was constructed specifically to be
harder for a small model: questions demanding a precise named subunit, a
specific consensus sequence, or a specific pair of interacting enzymes, of
the kind retrieval was expected to supply and memory was expected to
approximate. Ten control questions on core textbook material were included
to catch a miscalibrated set. The unaided model scored \textbf{0.955} on
the twenty-eight deliberately hard questions and 0.975 on the ten
controls, naming 130 of 136 required concepts without any retrieval at
all.

\textbf{Retrieval cannot beat memory at recalling material the model
already knows.} Its genuine value lies elsewhere, and Sections~\ref{4.4}
through~\ref{4.8} identify where.

\subsection*{A Limitation of This Finding}

The high baseline capability reported above is measured against this
project's own corpus, and that corpus draws heavily on public,
widely-available sources: Wikipedia, PubMed abstracts, open textbooks,
Gene Ontology, and UniProt. These are exactly the kind of material
commonly included in the training data of large language models. It is
therefore possible that the raw model answers most of this benchmark
correctly not because general biology knowledge is inherently easy for a
model this size, but because a meaningful share of the corpus, or material
closely resembling it, was already present in the model's own training
data. This would also explain why retrieval contributes so little: there
is little genuinely new information for retrieval to surface if the model
has already encountered the same public sources during pretraining.

The 38-question hard bank above is suggestive evidence against the
simplest version of this concern, since obscurity alone did not defeat the
model, but it does not rule out memorisation of the specific facts asked
about: a fact can be rarely discussed in ordinary text while still
appearing, close to verbatim, in a public source a pretraining run would
have scraped. Separating ``the model has learned biology'' from ``the
model has memorised this corpus'' would require a dedicated contamination
probe, for instance testing whether the model can continue specific
corpus passages verbatim, which is outside the scope of the measurements
in this chapter. It is recorded here because it directly affects how the
baseline-capability finding above should be read, and because a bare
accuracy number cannot distinguish between the two explanations on its
own.

\section{Hallucination Rate}\label{4.4}

Adversarial questions presuppose an entity that does not exist. Reading
the eight adversarial answers directly, rather than relying only on the
aggregate score:

\begin{table}[h]
\centering
\caption{Hallucination on the 8 adversarial questions}
\label{tab:4.3}
\begin{tabular}{@{}lr@{}}
\toprule
\textbf{System} & \textbf{Fabricates} \\
\midrule
Baseline & \textbf{1 / 8} \\
Full pipeline & \textbf{1 / 8} \\
\bottomrule
\end{tabular}
\end{table}

Both systems correctly refuse the false premise \textbf{7 times out of
8}, and the one genuine failure differs between them: the baseline invents
a mechanism for a fabricated ``law'', the full pipeline invents a
fabricated cycle name for a different question.

\section{Verification Layer: Calibration Against Real Output}\label{4.6}

\subsection{Deterministic Checkers}\label{4.6.2}

Every deterministic checker (Table~\ref{tab:3.2}) is calibrated against
537 real, previously generated answers, not synthetic test cases, since a
checker that only ever sees synthetic examples can pass every test it was
given and still misfire on real model output. Across this set the
deterministic layer raises \textbf{zero false blocking errors}.

\subsection{Contradiction Detection}\label{4.6.3}

Claim support is labelled \texttt{supported}, \texttt{contradictory}, or
\texttt{unsupported} against the retrieved evidence using a
natural-language-inference model, measured on twelve labelled test cases
already used to evaluate the external critic (Section~\ref{4.7}):

\begin{table}[h]
\centering
\caption[Contradiction detection performance]{Contradiction detection
performance. False alarms are measured on answers known to be sound.}
\label{tab:4.4}
\begin{tabular}{@{}lcc@{}}
\toprule
\textbf{Method} & \textbf{Contradictions caught} & \textbf{False alarms} \\
\midrule
NLI entailment & \textbf{5 / 7} & \textbf{0 / 5} \\
\bottomrule
\end{tabular}
\end{table}

The two contradictions not caught are, on inspection, correctly left to
other checkers: a reversed process order is caught by the process-order
checker, and a level mismatch is caught by the external critic, neither of
which is a factual contradiction in the sense the NLI check is built to
detect. Blocking is gated on what \emph{share} of an answer's claims are
contradicted, not on whether any single one is, since a real teaching
answer carries thirty or more claims and a single contradicted claim among
them is a warning worth recording rather than grounds to force a rewrite.

\section{External Critic Evaluation}\label{4.7}

An external critic model (GLM \texttt{glm-4.7-flash}) reviews generated
answers as an independent second opinion, bound to never override a
deterministic verifier error and to only ever make the pipeline's decision
\emph{stricter}, never looser (Section~\ref{3.7}). Measured on twelve
labelled cases with planted errors:

\begin{table}[h]
\centering
\caption{External critic versus the deterministic verifier alone}
\label{tab:4.5}
\begin{tabular}{@{}lccc@{}}
\toprule
\textbf{Method} & \textbf{Decision acc.} & \textbf{False accepts} & \textbf{False revises} \\
\midrule
Deterministic verifier & 0.750 & -- & -- \\
\textbf{GLM critic} & \textbf{0.917} & 1 / 12 & \textbf{0 / 12} \\
\bottomrule
\end{tabular}
\end{table}

The critic's one error is a false accept, the more dangerous of the two
possible mistakes since it lets a bad answer through rather than
needlessly flagging a good one, and it occurs on the case involving a
wholly fabricated organelle, which the deterministic entity-authority
check (Section~\ref{4.8}) independently catches from the other side. Zero
false-revises across twelve cases means the critic never penalised a sound
answer.

\section{Entity Authority Check}\label{4.8}

A lexicon of 135{,}941 entity names is built from the corpus's own naming
authorities (Gene Ontology, MeSH, Reactome, and the encyclopaedia index)
to answer a narrow but useful question: does this corpus, across several
independent sources, recognise this name at all? Applied to the
\emph{question} rather than the answer, this check identifies all eight
fabricated entities in the adversarial benchmark with a false-alarm rate
of 0.7\% (3 of 433) on genuine questions. It is applied only on the
question side, where it can add a cautionary note to the prompt but never
block a real answer outright. The same lexicon applied to an answer was
measured to be far less precise.

\section{Retrieval and Generation Performance}\label{4.9}

Mean retrieval latency is under one second once the index is warmed into
the operating system's page cache, with the worst-case query measured at
0.37 seconds. The local model runs fully GPU-offloaded at approximately
19.6 tokens per second on the project's available hardware. A fair
ablation requires every system under comparison to run under identical
generation settings (response-length budget, GPU offload), which this
project's harness enforces.

\section{Learned Question Router}\label{4.10}

Every question is routed on five dimensions before retrieval: task type,
biology subdomain, learner level, teaching style, and context depth
(Section~\ref{3.4}). Three of these dimensions, task type, subdomain, and
learner level, are additionally served by a learned classifier: frozen
sentence-embedding features (a general-purpose encoder and a biomedical
one, \cite{xiao2023bge,liu2021sapbert}) combined with lexical features,
feeding a linear or kernel classifier selected by cross-validated grid
search on macro F1, with questions in the held-out evaluation set strictly
excluded from training. The classifier is consulted alongside a
rule-based router and is permitted to override it only on dimensions
where it is measured to perform clearly better, so a hand-written rule
remains authoritative wherever it is not.

\begin{table}[h]
\centering
\caption[Rule-based versus learned routing]{Rule-based versus learned
routing on the held-out set (macro F1, fifty questions never seen in
training)}
\label{tab:4.5b}
\begin{tabular}{lccc}
\hline
\textbf{Dimension} & \textbf{Classes} & \textbf{Rules} & \textbf{Learned} \\
\hline
Learner level & 6 & 0.030 & \textbf{0.412} \\
Biology subdomain & 18 & 0.417 & \textbf{0.699} \\
Task type & 8 & 0.656 & \textbf{0.711} \\
\hline
\end{tabular}
\end{table}

The learned router outperforms the rules on every dimension, most sharply
on learner level, which selects the prompt's register: a postgraduate
question answered at the wrong level is answered in language reserved for
school students, which plausibly costs exactly the terminology the answer
keys reward.

\section{The Ablation Comparison}\label{4.11}

The core question this project set out to answer is what each layer of
the pipeline actually contributes on top of the raw model, with every
confounding difference in generation settings removed. All four systems
were run on the same held-out question set under identical generation
settings.

\begin{table}[h]
\centering
\caption[Four-system ablation, matched settings]{Four-system ablation,
matched settings, held-out question set. Latency is the median per answer.}
\label{tab:4.6}
\small
\begin{tabular}{@{}lccc@{}}
\toprule
\textbf{System} & \textbf{Overall} & \textbf{Concept recall} & \textbf{Latency} \\
\midrule
Baseline (raw model) & 0.887 & 0.894 & 35.4\,s \\
+ Hybrid retrieval only & 0.881 & 0.891 & 64.4\,s \\
+ Rerank, verify, revise (full pipeline) & 0.889 & 0.897 & 139.6\,s \\
+ External critic & 0.892 & 0.900 & 110.4\,s \\
\bottomrule
\end{tabular}
\end{table}

On raw accuracy the four systems are close enough to be indistinguishable
at this sample size, which is expected: Section~\ref{4.3} already
establishes that the raw model's own memory covers the large majority of
what this benchmark asks for, leaving little room for any addition to
move the accuracy needle. That is not where this pipeline earns its
place.

\textbf{Every claim the full pipeline makes is traceable; the baseline's
are not, at any accuracy.} Measured across the same fifty answers, 60.5\%
of the full pipeline's factual claims carry an inline citation, and of
those citations 65.2\% point at a passage that plausibly supports the
claim. The raw model scores zero on both, by construction, since it names
no source for anything it states. This is the capability an accuracy
score cannot show and the baseline cannot produce at all: a natural-
language-inference model that catches contradictions a word-overlap
checker structurally could not (Section~\ref{4.6.3}), an external critic
that catches errors the deterministic layer alone misses
(Section~\ref{4.7}), and a corpus-derived entity lexicon that catches
fabricated premises with a 0.7\% false-alarm rate (Section~\ref{4.8}).
None of those depend on beating the raw model at recalling biology from
memory, and all of them are things the raw model cannot do for itself at
any accuracy.

\section{Discussion}\label{4.12}

Two conclusions follow directly from the measurements in this chapter.

First, \textbf{the underlying small model was never the weak link}. It
already answers the large majority of this benchmark's questions
correctly unaided, and it already refuses a fabricated premise correctly
seven times out of eight (Section~\ref{4.4}). Retrieval-augmented
generation, built the way most published systems build it, cannot improve
on that by much, because its unique contribution is bounded by how much
the question set asks for that the model does not already know
(Section~\ref{4.3}), and that bound is narrow here.

Second, and more encouragingly, \textbf{the components that were built to
check the model's work each independently earn their place}, measured on
their own terms rather than assumed.

\begin{itemize}
\item Deterministic recomputation catches real arithmetic and factual
errors with a demonstrated near-zero false-positive rate on real output
(Section~\ref{4.6.2}).
\item A natural-language-inference model catches contradictions that a
simpler overlap heuristic structurally could not (Section~\ref{4.6.3}).
\item An external critic model adds real error-catching capacity beyond
the deterministic layer (Section~\ref{4.7}).
\item A corpus-derived entity lexicon recognises a fabricated premise with
high precision purely from local data (Section~\ref{4.8}).
\end{itemize}

None of these depend on out-scoring the raw model at answering biology
questions from memory. They depend on being able to show, for any
specific claim in any specific answer, whether it can be traced to a
source, whether it agrees with that source, and whether its arithmetic is
actually correct, which the raw model, however good its memory, cannot do
for itself at all. Chapter~5 argues that this, and not a marginal accuracy
gain, is the honest basis on which this project's contribution should be
judged.
